\documentclass{article}
\usepackage[utf8]{inputenc}
\usepackage{amsmath}
\usepackage{amssymb}
\usepackage{hyperref}
\usepackage{html}

\setcounter{secnumdepth}{0}
\setlength{\parindent}{0pt}
\setlength{\parskip}{0.75em}
\providecommand{\tightlist}{}

\title{Making Confer.to Work on Linux with TPM-backed WebAuthn}
\author{Vitor Py}
\date{2025-12-25}

\begin{document}
\begin{rawhtml}
<nav aria-label="Main navigation">
  <a href="/">Vitor Py</a> &middot;
  <a href="/blog/">Random Notes</a>
</nav>
<hr>
<script async src="https://plausible.io/js/pa-VS2SSXt2oRqjFWIl0ywIH.js"></script>
<script>
  window.plausible=window.plausible||function(){(plausible.q=plausible.q||[]).push(arguments)},plausible.init=plausible.init||function(i){plausible.o=i||{}};
  plausible.init()
</script>
\end{rawhtml}

\maketitle

\href{https://confer.to}{Confer.to} is an end-to-end encrypted AI chat service. Unlike typical AI assistants where your conversations feed into corporate training data, Confer uses the WebAuthn PRF extension to derive encryption keys directly from your fingerprint. No passwords, no server-side key storage---just your biometric tied to a hardware-backed credential.

There's a catch: this requires a \textbf{platform authenticator} with PRF support. Windows Hello and macOS Touch ID work out of the box. Linux? Not so much.

\section{The Problem: Linux Has No Platform Authenticator}

WebAuthn distinguishes between two types of authenticators:

\begin{itemize}
\item \textbf{Platform authenticators}: Built into your device (Windows Hello, Touch ID, Android biometrics)

\item \textbf{Roaming authenticators}: External hardware like YubiKeys

\end{itemize}

Linux has excellent support for roaming authenticators via USB/NFC. But platform authenticators? The browser expects something that simply doesn't exist in the standard Linux stack.

You might think "I'll just use my YubiKey." Unfortunately, most hardware security keys don't support the PRF extension. Even the YubiKey 5 series, which supports CTAP2.1 and \texttt{hmac-secret}, doesn't expose PRF to the browser in a way that works for client-side encryption use cases.

\section{The PRF Extension: Deriving Keys from Credentials}

The PRF (Pseudo-Random Function) extension solves a specific problem: how do you derive cryptographic material from a WebAuthn credential without storing secrets anywhere?

Here's how it works:

\begin{enumerate}
\item Website provides a \textbf{salt} (32 bytes)

\item Authenticator uses credential-specific secret material

\item Authenticator computes \texttt{HMAC-SHA256(credRandom, salt)} → 32-byte output

\item Website uses this output as an encryption key

\end{enumerate}

The magic: \textbf{same credential + same salt = same key, every time}. No password to remember, no key to store. Your fingerprint \emph{is} the key derivation mechanism.

\section{The Solution: TPM-FIDO2-PRF}

I built a two-component system that creates a platform authenticator for Linux:

\begin{itemize}
\item \textbf{\href{https://github.com/vitorpy/tpm-fido2-prf}{tpm-fido2-prf}}: A native Go application that implements FIDO2/CTAP2 using your system's TPM

\item \textbf{\href{https://github.com/vitorpy/tpm-fido2-extension}{tpm-fido2-extension}}: A Chrome extension that intercepts WebAuthn API calls

\end{itemize}

Together, they make the browser believe it has a built-in platform authenticator with full PRF support.

\subsection{Architecture Overview}

\textbf{Text fallback} (JavaScript disabled):

\begin{verbatim}
Website → inject.js → content.js → background.js → tpm-fido → fprintd → TPM
                                                              ↓
Website ← inject.js ← content.js ← background.js ← tpm-fido ← PRF + attestation
\end{verbatim}

\section{Deep Dive: How It Works}

\subsection{Chrome Extension Architecture}

The extension uses three layers to intercept WebAuthn:

\begin{enumerate}
\item \textbf{inject.js} (page context): Overrides \texttt{navigator.credentials.create()} and \texttt{navigator.credentials.get()}. Serializes \texttt{ArrayBuffer} objects to Base64 for JSON transport.

\item \textbf{content.js} (isolated world): Bridges between page context and extension via \texttt{window.postMessage} and \texttt{chrome.runtime.sendMessage}.

\item \textbf{background.js} (service worker): Communicates with the native host via Chrome's Native Messaging API.

\end{enumerate}

The extension also overrides \texttt{PublicKeyCredential.isUserVerifyingPlatformAuthenticatorAvailable()} to return \texttt{true}, making websites believe a platform authenticator exists.

\subsection{TPM Key Derivation}

This is the clever part. The TPM backend uses HKDF to derive unique keys per relying party:

\begin{verbatim}
// From tpm/tpm.go
info := append([]byte("tpm-fido-application-key"), applicationParam...)
r := hkdf.New(sha256.New, seed, []byte{}, info)
\end{verbatim}

On registration:

\begin{enumerate}
\item Generate a random 20-byte seed

\item Derive a primary key template using HKDF with the seed + RP ID hash

\item Create a TPM primary key under the Owner hierarchy

\item Generate a signing child key from the primary

\item Return handle: \texttt{{[}child\_private\_blob \textbar{} child\_public\_blob \textbar{} seed{]}}

\end{enumerate}

On authentication:

\begin{enumerate}
\item Extract seed from the handle

\item Reconstruct the primary key using the same HKDF derivation

\item Load the child key from the handle

\item Sign the challenge

\end{enumerate}

\textbf{Why this works}: The TPM cryptographically binds the child key to its parent. If the seed is wrong, the parent reconstruction fails, and the child key can't be loaded. Keys can't be extracted or used on a different TPM.

\subsection{PRF Implementation}

During credential creation, PRF outputs are computed immediately:

\begin{verbatim}
// Salt hashing per WebAuthn spec
hashedSalt := sha256.Sum256(append([]byte("WebAuthn PRF\x00"), salt...))

// Derive credential-specific random
credRandom := hmac.New(sha256.New, seed)
credRandom.Write([]byte("credential-random"))

// Compute PRF output
output := hmac.New(sha256.New, credRandom.Sum(nil))
output.Write(hashedSalt[:])
// output.Sum(nil) is the 32-byte PRF result
\end{verbatim}

The website receives this output and can use it directly as an AES key for encryption.

\subsection{User Verification via Fingerprint}

Every operation requires fingerprint verification:

\begin{verbatim}
// Spawn fprintd-verify and wait for result
cmd := exec.Command("fprintd-verify")
// ... handle result
\end{verbatim}

A desktop notification prompts the user to touch the fingerprint reader. This satisfies the "user verification" requirement of WebAuthn.

\section{Setup Guide}

\subsection{Prerequisites}

\begin{enumerate}
\item \textbf{TPM 2.0}: Check with \texttt{ls /dev/tpmrm0}

\item \textbf{TPM access}: Add yourself to the \texttt{tss} group:

\begin{verbatim}
sudo usermod -aG tss $USER
# Log out and back in
\end{verbatim}

\item \textbf{Fingerprint reader}: Must be supported by fprintd

\item \textbf{Enrolled fingerprint}:

\begin{verbatim}
fprintd-enroll
\end{verbatim}

\item \textbf{Go 1.21+}: For building the native component

\item \textbf{Chrome or Chromium}

\end{enumerate}

\subsection{Clone and Build}

\begin{verbatim}
# Clone both repos
git clone https://github.com/vitorpy/tpm-fido2-prf.git
git clone https://github.com/vitorpy/tpm-fido2-extension.git

# Build native component
cd tpm-fido2-prf
go build -o ~/bin/tpm-fido .
\end{verbatim}

\subsection{Load the Chrome Extension}

\begin{enumerate}
\item Open Chrome → \texttt{chrome://extensions}

\item Enable "Developer mode"

\item Click "Load unpacked" → select the \texttt{tpm-fido2-extension} directory

\item The extension ID should be \texttt{bfmfknknibchmioeamgbnlpakcjimnbf} (stable, derived from manifest key)

\end{enumerate}

\subsection{Configure Native Messaging}

\begin{verbatim}
cd tpm-fido2-extension
./setup.sh
\end{verbatim}

This creates \texttt{\textasciitilde{}/.config/google-chrome/NativeMessagingHosts/com.vitorpy.tpmfido.json}:

\begin{verbatim}
{
  "name": "com.vitorpy.tpmfido",
  "description": "TPM-FIDO WebAuthn Platform Authenticator with PRF support",
  "path": "/home/youruser/bin/tpm-fido",
  "type": "stdio",
  "allowed_origins": ["chrome-extension://bfmfknknibchmioeamgbnlpakcjimnbf/"]
}
\end{verbatim}

\subsection{Test the Setup}

\begin{enumerate}
\item Restart Chrome

\item Visit \href{https://webauthn.io}{webauthn.io}

\item Register a credential → fingerprint prompt should appear

\item Authenticate → verify it works

\end{enumerate}

\textbf{Debug}: \texttt{chrome://extensions} → TPM-FIDO2 → "Inspect views: service worker"

\subsection{Using with Confer.to}

Once set up, navigate to \href{https://confer.to}{confer.to} and login. The platform authenticator will be used automatically.

\section{Security Considerations}

\begin{itemize}
\item \textbf{Private keys never leave the TPM}: Only the TPM can sign; key material is hardware-protected

\item \textbf{Fingerprint required for every operation}: No silent authentication

\item \textbf{Credential metadata stored locally}: Only user info and credential IDs at \texttt{\textasciitilde{}/.local/share/tpm-fido/credentials.json}---not keys

\item \textbf{Self-signed attestation}: Uses a generic certificate to avoid device fingerprinting (privacy trade-off vs. strict RP requirements)

\end{itemize}

\section{Limitations}

\begin{itemize}
\item \textbf{Chrome/Chromium only}: No Firefox WebExtension port (yet)

\item \textbf{Requires fprintd}: No alternative user verification methods

\item \textbf{Self-signed attestation}: Some strict relying parties may reject credentials without hardware attestation

\end{itemize}

\section{Repository Links}

\begin{itemize}
\item \href{https://github.com/vitorpy/tpm-fido2-prf}{tpm-fido2-prf} - Native TPM-backed authenticator

\item \href{https://github.com/vitorpy/tpm-fido2-extension}{tpm-fido2-extension} - Chrome extension

\item \href{https://confer.to}{Confer.to} - E2E encrypted AI chat

\item \href{https://confer.to/blog/2025/12/confessions-to-a-data-lake/}{Confessions to a Data Lake} - Confer's privacy philosophy

\end{itemize}

\begin{rawhtml}
<hr>
<footer>
  <a rel="license" href="https://creativecommons.org/licenses/by/4.0/">This work is licensed under a Creative Commons Attribution 4.0 International License</a>.<br>
  Built with <a href="/build/">LaTeX2HTML</a>.
</footer>
\end{rawhtml}

\end{document}
